%% main.tex — generated by cn-latex-kit
%% Compile with XeLaTeX (or just run your normal build; .latexmkrc handles it).
\documentclass[UTF8]{journaltemplate}

\cnTitle{哈尔滨理工大学学报}
\cnAuthor{作者一$^{1}$，作者二$^{2}$}
\cnAffil{(1. 单位一，城市 邮编；2. 单位二，城市 邮编)}
\cnAbstract{摘要内容：按报道性摘要撰写，给出目的、过程、结果和结论，字数一般 250--300 汉字。}
\cnKeywords{关键词一；关键词二；关键词三}
\doi{10.15938/j.example.2016.01.001}
\clc{TB 114}
\docCode{A}
\articleID{1007-449X(2016)01-0000-00}
\enTitle{English Title of the Paper}
\enAuthor{AUTHOR One$^{1}$, AUTHOR Two$^{2}$}
\enAffil{(1. Affiliation One, City, Country; 2. Affiliation Two, City, Country)}
\enAbstract{Contents of the abstract in English.}
\enKeywords{keyword one; keyword two; keyword three}
\received{2016--01--01}
\foundation{国家自然科学基金（编号）}
\authorBio{作者姓名（出生年—），性别，学位，职称，研究方向。}
\corresponding{作者姓名}
\runningAuthor{作者等}
\runningTitle{哈尔滨理工大学学报}
\journalName{哈尔滨理工大学学报}
\journalNameEN{Journal Template}
\journalVol{21}
\journalIssue{1}
\journalYear{2016}
\journalMonth{2}
\journalMonthEN{Feb.}

\begin{document}
\twocolumn[\maketitle]

\section{引\;言}
在此撰写引言：阐述研究背景、存在问题、本文探讨范围与采用的方法。
文中首次出现的缩写应给出全称，如开关磁阻电机（switched reluctance motor, SRM）。
参考文献在正文中以上标标出，如\upcite{1}。

\section{正文标题}
在此撰写正文。正文首行缩进 2 字，字体字号与行距已由模板自动设置。

\subsection{小节标题}
正文中的图、表、公式一律用阿拉伯数字编号，如图~\ref{fig:sample}、式~\ref{eq:sample}。

\begin{figure}[H]
  \centering
  \includegraphics[width=0.72\columnwidth]{example-image}
  \caption{示例图名}
  \figCaptionEN{An example figure caption}
  \label{fig:sample}
\end{figure}

\begin{table}[H]
  \centering
  \caption{示例表名}
  \tabCaptionEN{An example table caption}
  \label{tab:sample}
  \begin{tabular}{lc}
    \toprule
    参数 & 数值 \\
    \midrule
    直流电源 $V_{dc}$/V & 100 \\
    电阻 $R/\Omega$     & 0.09 \\
    \bottomrule
  \end{tabular}
\end{table}

定义公式形如
\begin{equation}
  T_e = \frac{3}{2}p\left[\psi_f i_q + (L_d - L_q)i_d i_q\right].
  \label{eq:sample}
\end{equation}
式中：$T_e$ 为电磁转矩；$i_d$、$i_q$ 为 $d$、$q$ 轴电流。

\section{结\;论}
在此撰写结论，给出详细、充分的结论性语言。

\begin{thebibliography}{99}
\bibitem{1} 作者. 文献题名[J]. 刊名, 年, 卷(期): 页码.
\bibitem{2} AUTHOR A B. Article title[J]. Journal Name, 2016, 21(1): 1--10.
\end{thebibliography}

\end{document}
